% ============================================================
% CATT-CCS: Domain-Constrained Adversarial Evaluation of NIDS
%           Across Classifier Architectures
% Target: ACM CCS 2027 — Security and Privacy of Machine Learning track
% Extends: AISec '26 workshop paper (CATT)
% ============================================================
% Compile: pdflatex catt_ccs.tex → bibtex catt_ccs → pdflatex × 2
% ============================================================

\documentclass[sigconf,anonymous=true]{acmart}

\usepackage{booktabs}
\usepackage{multirow}
\usepackage{xcolor}
\usepackage{listings}
\usepackage{microtype}
\usepackage{amsmath}
\usepackage{algorithm}
\usepackage{algorithmic}

\acmConference[CCS '27]{ACM Conference on Computer and Communications Security}{2027}{TBD}
\acmDOI{}
\acmISBN{}
\copyrightyear{2027}
\acmYear{2027}
\setcopyright{acmlicensed}

\definecolor{deltagreen}{RGB}{0,140,0}
\definecolor{deltared}{RGB}{180,0,0}

\newcommand{\gap}[1]{\textcolor{deltagreen}{+#1}}
\newcommand{\con}{\text{con}}
\newcommand{\uncon}{\text{uncon}}

% =============================================================
\begin{document}

\title{Constraint Inflation in Adversarial NIDS Evaluation:\\
       A Multi-Architecture, Multi-Dataset Analysis}

\author{[Author redacted for review]}
\affiliation{\institution{[Institution redacted for review]}}
\email{[redacted]}

% =============================================================
\begin{abstract}
Adversarial evasion evaluations of Network Intrusion Detection Systems
(NIDS) consistently overstate attack success because they omit domain
constraint projection from the inner loop of gradient-based attacks.
Without projection, optimized perturbations push features into physically
impossible regions---negative durations, sub-zero TTL values, rates
outside $[0,1]$---that cannot appear in observed network traffic.
Our prior work (CATT,~\cite{sasso2026catt}) quantified this
\emph{constraint inflation gap} in white-box settings against a single
MLP architecture on two datasets.

In this paper we ask: \emph{does the gap persist when attacks transfer
to the classifiers practitioners actually deploy?}
We evaluate constrained and unconstrained FGSM and PGD against three
classifier families---MLP (white-box surrogate), Random Forest, and
XGBoost---across UNSW-NB15, CICIDS-2017, and NSL-KDD with three
independent random seeds each.

We find that the gap is large and reproducible on tightly-constrained
datasets: \textbf{72.0 pp} on UNSW-NB15 ($\varepsilon=0.20$, PGD) and
\textbf{67.4 pp} on NSL-KDD ($\varepsilon=0.10$, PGD).
On CICIDS-2017, where most features carry only a non-negativity bound,
the gap is smaller (\textbf{12.2 pp}), confirming that inflation scales
with constraint tightness rather than being a dataset-specific artefact.
Transfer results are architecture- and dataset-dependent: MLP surrogates
transfer near-perfectly to RF and XGBoost on CICIDS-2017, but only
modestly on UNSW-NB15 and NSL-KDD, consistent with the looser feature
manifold geometry of the former.
We further show that constrained PGD is the \emph{optimal adaptive
adversary} within the $L_\infty +$ ConstraintSpec threat model, so our
constrained evasion rates are upper bounds on what a domain-aware
attacker can achieve.

The full benchmark suite is released at [anonymous repository].
\end{abstract}

\keywords{adversarial machine learning, network intrusion detection,
          domain constraints, transfer attacks, FGSM, PGD,
          UNSW-NB15, CICIDS-2017, NSL-KDD}

\maketitle

% =============================================================
\section{Introduction}
\label{sec:intro}

Network Intrusion Detection Systems are a primary defense layer in
enterprise and cloud environments, and adversarial robustness evaluations
of NIDS classifiers have proliferated in recent years~\cite{
  apruzzese2022modeling, han2021evaluating, yang2022carefully}.
The standard evaluation protocol---run FGSM or PGD against the trained
classifier, report evasion rate---has a fundamental flaw: gradient-based
attacks operate in unconstrained feature space and routinely produce
adversarial examples that are physically impossible to instantiate as
network traffic.

Consider a PGD attack that shifts the \texttt{sttl} (source
time-to-live) feature from its normalized representation of $0.8$ to
$-3.2$.  In original units this corresponds to a TTL of roughly $-119$,
an unsigned 8-bit field.  No network stack can generate such a packet;
no sensor can record it.  The evasion rate attributed to that perturbation
is not a property of the threat: it is an artefact of an unconstrained
optimizer exploiting infeasible regions of the feature space.

Our prior work~\cite{sasso2026catt} introduced CATT---a
\emph{ConstraintSpec} projection layer that clips each feature to its
documented domain bounds inside the PGD inner loop---and measured the
resulting \emph{constraint inflation gap} on UNSW-NB15 and CICIDS-2017
in a white-box MLP setting.  The gap was substantial: unconstrained PGD
overstated evasion by tens of percentage points at moderate $\varepsilon$.

This paper asks whether the finding generalizes along two axes that
CATT did not explore.

\noindent\textbf{Generalization across classifier architectures.}
Real NIDS deployments overwhelmingly use tree-based
classifiers~\cite{he2023survey}.  If the constraint inflation gap
collapses when adversarial examples transfer from an MLP surrogate to
RF or XGBoost, the AISec result is an MLP-specific curiosity rather than
a systemic evaluation flaw.

\noindent\textbf{Generalization across datasets.}
Different NIDS benchmark datasets encode different feature semantics
and constraint tightness.  UNSW-NB15 contains TTL fields, boolean
flags, and rate features with tight documented bounds; CICIDS-2017 uses
CICFlowMeter output where most features carry only a non-negativity
constraint; NSL-KDD has rate features strictly in $[0,1]$ and count
features with documented ceilings.  If the gap varies predictably with
constraint tightness, the mechanism is well-understood; if it does not,
the finding is fragile.

\noindent\textbf{Summary of findings.}
\begin{enumerate}
  \item The constraint inflation gap is large and statistically robust
        across three seeds on tightly-constrained datasets (72.0 pp,
        UNSW-NB15; 67.4 pp, NSL-KDD).
  \item The gap scales with constraint tightness: CICIDS-2017, with
        looser feature semantics, shows a smaller but consistently
        positive gap (up to 12.2 pp).
  \item Transfer results are architecture- and dataset-dependent.
        On CICIDS-2017, gradient-based examples from an MLP surrogate
        transfer at $\approx 100\%$ to RF and XGBoost regardless of
        whether constraints are applied.  On UNSW-NB15 and NSL-KDD,
        transfer rates are moderate (10--70\%) and highly variable
        across seeds.
  \item Constrained PGD is the optimal adaptive adversary against the
        ConstraintSpec projection; our constrained evasion rates are
        therefore upper bounds, not lower bounds.
\end{enumerate}

% =============================================================
\section{Threat Model}
\label{sec:threat}

We consider an adversary who seeks to evade a deployed NIDS classifier
by crafting malicious network flows that are misclassified as benign.
We formalize the threat model along four dimensions.

\subsection{Attacker Capability}

\textbf{White-box (surrogate).}
The adversary has full access to a surrogate model---weights,
architecture, and preprocessing pipeline---trained on a representative
subset of the same traffic distribution as the target.  This models an
adversary who has reverse-engineered or independently trained a
functionally equivalent model.

\textbf{Black-box (transfer).}
The adversary generates adversarial examples against the surrogate and
submits them to an unknown target classifier, receiving only binary
(attack/benign) decisions.  This is the realistic model for adversaries
who do not have access to deployed classifier internals.

\subsection{Threat Surface}

Attacks perturb the \emph{preprocessed feature vector} presented to the
classifier: the 40-dimensional StandardScaler-normalized representation
for UNSW-NB15 (plus one-hot encoded categorical features), the
78-dimensional all-numeric representation for CICIDS-2017, and the
38-dimensional representation for NSL-KDD (plus one-hot encoded
categorical features).
Constraint bounds are computed in the transformed space as
$\ell_i = (\text{lo}_i - \mu_i)/\sigma_i$ and
$u_i = (\text{hi}_i - \mu_i)/\sigma_i$.
Perturbations to categorical (one-hot) features are excluded; only
numeric features are perturbed, reflecting that categorical fields such
as protocol type are generally immutable at the network level.

We note that this threat model operates at the feature-vector level.
Translating feature-space perturbations back into actual packet sequences
is an open problem in adversarial NIDS research and out of scope here;
we follow the standard evaluation protocol of the
field~\cite{apruzzese2022modeling}.

\subsection{Practicality}

Perturbation budgets are expressed in both normalized and
original-domain units (Table~\ref{tab:calibration}).
At $\varepsilon = 0.20$ on UNSW-NB15, the budget corresponds to
approximately 0.70 seconds in connection duration, 14.9 source
packets, and 10,489 source bytes---perturbations achievable through
selective traffic shaping.
The same $\varepsilon$ on NSL-KDD corresponds to 496 additional
seconds of duration and 1.15~MB of source bytes: a much larger
original-domain budget, reflecting the heavy-tailed byte distributions
in that dataset.
This cross-dataset comparison motivates reporting $\varepsilon$ in
original units alongside the normalized value.

\subsection{Adaptive Adversary}
\label{sec:adaptive}

The strongest adaptive adversary against our ConstraintSpec projection
is one who knows the projection and crafts examples that survive
clamping.  We show in Appendix~\ref{app:adaptive} that
\emph{constrained PGD is precisely this adaptive attack}: at each step,
the adversary moves in the sign of the gradient and then projects onto
both the $\varepsilon$-ball and the valid feature manifold.  No strictly
stronger adaptive attack exists within the
$L_\infty + \text{ConstraintSpec}$ threat model.
Our constrained evasion rates are therefore \emph{upper bounds} on what
a domain-aware adversary can achieve.

% =============================================================
\section{Background and Related Work}
\label{sec:related}

\subsection{NIDS Feature Extraction and Semantics}

Modern NIDS classifiers operate on flow-level features extracted by
tools such as CICFlowMeter~\cite{lashkari2017cicflowmeter} or Argus.
These features encode hard physical constraints: packet lengths are
bounded by the MTU (typically 1,500 bytes), TTL is an 8-bit unsigned
field, inter-arrival times are non-negative, and flag counts cannot
exceed observed packet counts.  These constraints are
\emph{dataset-independent}: they follow from the TCP/IP specification
and are violated by any unconstrained gradient optimization that steps
outside the valid feature manifold.

\subsection{Adversarial Machine Learning}

FGSM~\cite{goodfellow2015explaining} computes a single-step
perturbation in the sign of the loss gradient:
$x' = x + \varepsilon \cdot \operatorname{sign}(\nabla_x \mathcal{L})$.
PGD~\cite{madry2018towards} iterates this step with step size $\alpha$,
projecting after each step onto the $L_\infty$ $\varepsilon$-ball around
the original input.  Both attacks are standard baselines in adversarial
ML evaluation.

\subsection{Transfer Attacks}

Papernot et al.~\cite{papernot2017practical} demonstrated that
adversarial examples generated against one model often transfer to a
different model trained on the same task, a property exploited in
black-box attacks.  Transfer success depends on the alignment of
decision boundaries between surrogate and target; gradient-based
surrogates transfer well to other differentiable models but less
reliably to piecewise-constant tree classifiers.

\subsection{Adaptive Attacks}

Tramer et al.~\cite{tramer2020adaptive} argue that security evaluations
must consider the strongest attacker who knows and exploits the
defense mechanism.  Our constrained PGD satisfies this requirement
for the ConstraintSpec projection (Section~\ref{sec:adaptive}).

\subsection{Adversarial Attacks on NIDS}

A growing body of work applies adversarial ML to NIDS
evaluation~\cite{han2021evaluating, yang2022carefully, lin2022idsgan,
debicha2023adv}.  These papers uniformly report evasion rates without
domain constraint projection, making their numbers incomparable to
physically achievable attack success.  Our work quantifies the resulting
overstatement.

% =============================================================
\section{Method: ConstraintSpec}
\label{sec:method}

\subsection{Formalization}

Let $\mathbf{x} \in \mathbb{R}^d$ be a preprocessed feature vector
(StandardScaler-normalized) presented to a NIDS classifier.  A
\emph{ConstraintSpec} $\mathcal{P}$ is a set of per-feature bounds
$\{[\ell_i, u_i]\}_{i=1}^{d}$ derived from the documented physical
semantics of the underlying feature, transformed into the preprocessed
space via $\ell_i = (\text{lo}_i - \mu_i)/\sigma_i$, where $\mu_i$ and
$\sigma_i$ are the fitted scaler statistics.

The \emph{ConstraintSpec projection} $\Pi_\mathcal{P}$ is the
element-wise clamp:
\[
  \left[\Pi_\mathcal{P}(\mathbf{x})\right]_i =
  \min\!\left(u_i,\, \max\!\left(\ell_i,\, x_i\right)\right).
\]
Features with no documented upper bound have $u_i = +\infty$; features
with no lower bound have $\ell_i = -\infty$.

\subsection{Inner-Loop Projection}

Standard PGD projects only onto the $L_\infty$ $\varepsilon$-ball after
each gradient step.  \emph{Constrained PGD} additionally applies
$\Pi_\mathcal{P}$:

\begin{algorithm}[h]
\caption{Constrained PGD}
\begin{algorithmic}[1]
\REQUIRE model $f_\theta$, input $\mathbf{x}_0$, label $y$,
         budget $\varepsilon$, step $\alpha$, steps $T$,
         spec $\mathcal{P}$
\STATE $\mathbf{x}^{(0)} \leftarrow \Pi_\mathcal{P}\!\left(
         \mathbf{x}_0 + \mathbf{u}\right)$
         where $\mathbf{u} \sim \mathcal{U}[-\varepsilon, \varepsilon]^d$
\FOR{$t = 1$ \TO $T$}
  \STATE $\mathbf{g} \leftarrow \nabla_{\mathbf{x}} \mathcal{L}(f_\theta,
         \mathbf{x}^{(t-1)}, y)$
  \STATE $\mathbf{x}^{(t)} \leftarrow \Pi_\varepsilon\!\left(
         \Pi_\mathcal{P}\!\left(
           \mathbf{x}^{(t-1)} + \alpha \cdot \operatorname{sign}(\mathbf{g})
         \right)\right)$
\ENDFOR
\RETURN $\mathbf{x}^{(T)}$
\end{algorithmic}
\end{algorithm}

\noindent
where $\Pi_\varepsilon(\mathbf{x}) =
\text{clip}(\mathbf{x},\, \mathbf{x}_0 - \varepsilon,\,
\mathbf{x}_0 + \varepsilon)$.
Setting $\ell_i = -\infty$, $u_i = +\infty$ for all $i$ recovers
standard (unconstrained) PGD.  The \emph{constraint inflation gap} is
the difference in evasion rate between unconstrained and constrained
attacks at the same $\varepsilon$.

\subsection{Dataset Constraint Specifications}

\textbf{UNSW-NB15.}  40 numeric features.  TTL fields are bounded to
$[0, 255]$; boolean indicators (\texttt{is\_sm\_ips\_ports}) to
$[0,1]$; all duration, packet count, and byte features to $[0, +\infty)$.

\textbf{CICIDS-2017.}  78 numeric features (CICFlowMeter output).
Flag counts (\texttt{FIN Flag Count}, etc.) are bounded to $[0, 255]$;
destination port to $[0, 65535]$; all remaining features to
$[0, +\infty)$.  Three inter-arrival time features
(\texttt{Flow IAT Min}, \texttt{Fwd IAT Min}, \texttt{Bwd IAT Min})
carry no lower bound---a known CICFlowMeter artefact where
out-of-order packet reordering produces negative IATs in observed
clean traffic.

\textbf{NSL-KDD.}  38 numeric features (KDD Cup feature set).
All 15 rate features are bounded to $[0, 1]$; \texttt{count} and
\texttt{srv\_count} to $[0, 511]$; \texttt{dst\_host\_count} and
\texttt{dst\_host\_srv\_count} to $[0, 255]$; 6 boolean features
to $[0, 1]$; remaining features to $[0, +\infty)$.

Each specification is validated via \texttt{validity\_report()} on
clean data before any attack is run; all three datasets produce zero
violations, confirming correct specification-to-pipeline alignment.

\subsection{Transfer Evaluation Protocol}

All adversarial examples are generated against the MLP surrogate.
Transfer evasion rate is then measured by running the same examples
through RF and XGBoost without retraining---the classifiers never see
the surrogate gradients.  This is the black-box transfer threat model
of~\cite{papernot2017practical}.

% =============================================================
\section{Empirical Evaluation}
\label{sec:eval}

\subsection{Experimental Setup}

\textbf{Datasets.}
UNSW-NB15~\cite{moustafa2015unsw}: 500,000 flows (stratified sample),
$\approx$87\% benign.
CICIDS-2017~\cite{sharafaldin2018generating}: 500,000 flows (stratified
sample) from the CICFlowMeter CSVs.
NSL-KDD~\cite{tavallaee2009detailed}: 147,341 flows (full KDDTrain+
plus KDDTest+), $\approx$53\% attack.

\textbf{Classifiers.}
\begin{itemize}
  \item \emph{MLP surrogate}: 3-layer (256--128--64), BatchNorm,
        dropout $p=0.3$, BCEWithLogitsLoss with imbalance weighting,
        Adam optimizer ($\text{lr}=10^{-3}$), early stopping on
        validation F1 (patience 5).
  \item \emph{Random Forest}: 300 trees, balanced class weights,
        \texttt{n\_jobs=-1}.
  \item \emph{XGBoost}: 400 trees, max depth 6, learning rate 0.1,
        \texttt{scale\_pos\_weight} tuned per dataset and seed.
\end{itemize}
All classifiers share the same StandardScaler preprocessing pipeline
fitted on the 80\% training split.

\textbf{Attack configuration.}  FGSM and PGD ($T=40$,
$\alpha = 2.5\varepsilon/40$, random start) with
$\varepsilon \in \{0.05, 0.10, 0.20, 0.30, 0.40, 0.50\}$.
Evasion rate is computed over attack-class samples only: the fraction
of true attacks the classifier misclassifies as benign after
perturbation.  Three random seeds per configuration; results reported
as mean $\pm$ standard deviation.

\subsection{White-Box MLP Results}

Tables~\ref{tab:unsw_wb}--\ref{tab:nsl_wb} report white-box MLP evasion
rates across all three datasets.  The headline finding is consistent:
unconstrained attacks overstate evasion on every dataset, at every
$\varepsilon$, in all 3 seeds.

\begin{table}[t]
\centering
\caption{UNSW-NB15 white-box MLP evasion rate (\%), mean $\pm$ std
         across 3 seeds.  $\Delta$ = unconstrained $-$ constrained.
         Peak gap at $\varepsilon=0.20$: \textbf{+72.0 pp} (PGD).}
\label{tab:unsw_wb}
\small
\begin{tabular}{r rr r rr r}
\toprule
 & \multicolumn{3}{c}{\textbf{FGSM}} & \multicolumn{3}{c}{\textbf{PGD}} \\
\cmidrule(lr){2-4}\cmidrule(lr){5-7}
$\varepsilon$ & Con. & Uncon. & $\Delta$ & Con. & Uncon. & $\Delta$ \\
\midrule
0.05 & $0.02{\scriptstyle\pm0.02}$ & $0.16{\scriptstyle\pm0.06}$ & \gap{0.14} &
       $0.02{\scriptstyle\pm0.02}$ & $0.20{\scriptstyle\pm0.06}$ & \gap{0.17} \\
0.10 & $0.17{\scriptstyle\pm0.13}$ & $2.51{\scriptstyle\pm1.28}$ & \gap{2.34} &
       $0.21{\scriptstyle\pm0.14}$ & $3.68{\scriptstyle\pm1.37}$ & \gap{3.47} \\
0.20 & $2.92{\scriptstyle\pm0.60}$ & $54.2{\scriptstyle\pm36.3}$ & \gap{51.3} &
       $7.14{\scriptstyle\pm3.69}$ & $79.1{\scriptstyle\pm19.4}$ & \gap{72.0} \\
0.30 & $21.5{\scriptstyle\pm5.16}$ & $81.9{\scriptstyle\pm28.0}$ & \gap{60.4} &
       $41.8{\scriptstyle\pm8.72}$ & $99.5{\scriptstyle\pm0.38}$ & \gap{57.7} \\
0.40 & $54.2{\scriptstyle\pm14.8}$ & $97.0{\scriptstyle\pm4.84}$ & \gap{42.8} &
       $97.8{\scriptstyle\pm0.44}$ & $99.5{\scriptstyle\pm0.54}$ & \gap{1.72} \\
0.50 & $83.8{\scriptstyle\pm20.7}$ & $99.4{\scriptstyle\pm0.90}$ & \gap{15.6} &
       $99.4{\scriptstyle\pm0.29}$ & $99.2{\scriptstyle\pm0.92}$ & $-0.16$ \\
\bottomrule
\end{tabular}
\end{table}

\begin{table}[t]
\centering
\caption{CICIDS-2017 white-box MLP evasion rate (\%), mean $\pm$ std
         across 3 seeds.  Peak gap at $\varepsilon=0.30$: \textbf{+12.2 pp} (PGD).
         Smaller gap reflects looser feature constraints (see
         Section~\ref{sec:discussion}).}
\label{tab:cicids_wb}
\small
\begin{tabular}{r rr r rr r}
\toprule
 & \multicolumn{3}{c}{\textbf{FGSM}} & \multicolumn{3}{c}{\textbf{PGD}} \\
\cmidrule(lr){2-4}\cmidrule(lr){5-7}
$\varepsilon$ & Con. & Uncon. & $\Delta$ & Con. & Uncon. & $\Delta$ \\
\midrule
0.05 & $16.7{\scriptstyle\pm4.51}$ & $24.4{\scriptstyle\pm4.33}$ & \gap{7.7} &
       $22.4{\scriptstyle\pm4.29}$ & $33.4{\scriptstyle\pm4.21}$ & \gap{11.0} \\
0.10 & $35.5{\scriptstyle\pm4.88}$ & $48.4{\scriptstyle\pm2.45}$ & \gap{12.9} &
       $44.5{\scriptstyle\pm2.82}$ & $54.8{\scriptstyle\pm0.88}$ & \gap{10.3} \\
0.20 & $56.4{\scriptstyle\pm1.11}$ & $62.0{\scriptstyle\pm1.15}$ & \gap{5.6} &
       $61.2{\scriptstyle\pm0.36}$ & $67.9{\scriptstyle\pm1.90}$ & \gap{6.8} \\
0.30 & $65.9{\scriptstyle\pm2.73}$ & $75.7{\scriptstyle\pm4.97}$ & \gap{9.8} &
       $74.9{\scriptstyle\pm3.56}$ & $87.1{\scriptstyle\pm3.82}$ & \gap{12.2} \\
0.40 & $81.7{\scriptstyle\pm6.45}$ & $89.3{\scriptstyle\pm5.06}$ & \gap{7.6} &
       $90.0{\scriptstyle\pm4.17}$ & $96.7{\scriptstyle\pm1.89}$ & \gap{6.6} \\
0.50 & $88.8{\scriptstyle\pm5.50}$ & $95.4{\scriptstyle\pm1.58}$ & \gap{6.6} &
       $96.7{\scriptstyle\pm2.39}$ & $99.3{\scriptstyle\pm0.32}$ & \gap{2.6} \\
\bottomrule
\end{tabular}
\end{table}

\begin{table}[t]
\centering
\caption{NSL-KDD white-box MLP evasion rate (\%), mean $\pm$ std
         across 3 seeds.  Peak gap at $\varepsilon=0.10$: \textbf{+67.4 pp} (PGD),
         reflecting tight rate-feature constraints.}
\label{tab:nsl_wb}
\small
\begin{tabular}{r rr r rr r}
\toprule
 & \multicolumn{3}{c}{\textbf{FGSM}} & \multicolumn{3}{c}{\textbf{PGD}} \\
\cmidrule(lr){2-4}\cmidrule(lr){5-7}
$\varepsilon$ & Con. & Uncon. & $\Delta$ & Con. & Uncon. & $\Delta$ \\
\midrule
0.05 & $4.62{\scriptstyle\pm0.62}$ & $10.98{\scriptstyle\pm0.85}$ & \gap{6.36} &
       $5.20{\scriptstyle\pm0.54}$ & $13.0{\scriptstyle\pm1.35}$ & \gap{7.85} \\
0.10 & $13.1{\scriptstyle\pm1.57}$ & $58.9{\scriptstyle\pm13.8}$ & \gap{45.8} &
       $16.7{\scriptstyle\pm2.08}$ & $84.1{\scriptstyle\pm6.13}$ & \gap{67.4} \\
0.20 & $78.7{\scriptstyle\pm4.60}$ & $95.7{\scriptstyle\pm2.55}$ & \gap{17.0} &
       $90.5{\scriptstyle\pm4.69}$ & $98.7{\scriptstyle\pm0.73}$ & \gap{8.20} \\
0.30 & $93.7{\scriptstyle\pm3.36}$ & $98.8{\scriptstyle\pm1.11}$ & \gap{5.1} &
       $97.7{\scriptstyle\pm1.19}$ & $99.6{\scriptstyle\pm0.37}$ & \gap{1.84} \\
0.40 & $97.1{\scriptstyle\pm1.69}$ & $99.4{\scriptstyle\pm0.69}$ & \gap{2.3} &
       $99.5{\scriptstyle\pm0.53}$ & $99.8{\scriptstyle\pm0.15}$ & \gap{0.36} \\
0.50 & $98.8{\scriptstyle\pm0.89}$ & $99.6{\scriptstyle\pm0.58}$ & \gap{0.8} &
       $99.6{\scriptstyle\pm0.43}$ & $99.9{\scriptstyle\pm0.06}$ & \gap{0.24} \\
\bottomrule
\end{tabular}
\end{table}

\subsection{Transfer Attack Results}
\label{sec:transfer}

Table~\ref{tab:transfer} reports transfer evasion rates for PGD
adversarial examples (generated against the MLP surrogate) evaluated on
RF and XGBoost across all three datasets at representative
$\varepsilon$ values.

\begin{table*}[t]
\centering
\caption{Transfer evasion rate (\%): PGD adversarial examples generated
         against MLP surrogate, evaluated on black-box RF and XGBoost.
         Mean $\pm$ std, 3 seeds.  Con.~= constrained PGD;
         Uncon.~= unconstrained PGD.}
\label{tab:transfer}
\small
\begin{tabular}{l r rr rr rr rr rr rr}
\toprule
 & & \multicolumn{4}{c}{\textbf{UNSW-NB15}} &
     \multicolumn{4}{c}{\textbf{CICIDS-2017}} &
     \multicolumn{4}{c}{\textbf{NSL-KDD}} \\
\cmidrule(lr){3-6}\cmidrule(lr){7-10}\cmidrule(lr){11-14}
$\varepsilon$ & & \multicolumn{2}{c}{RF} & \multicolumn{2}{c}{XGB} &
                   \multicolumn{2}{c}{RF} & \multicolumn{2}{c}{XGB} &
                   \multicolumn{2}{c}{RF} & \multicolumn{2}{c}{XGB} \\
\cmidrule(lr){3-4}\cmidrule(lr){5-6}\cmidrule(lr){7-8}\cmidrule(lr){9-10}
\cmidrule(lr){11-12}\cmidrule(lr){13-14}
 & & Con. & Uncon. & Con. & Uncon. & Con. & Uncon. & Con. & Uncon. &
     Con. & Uncon. & Con. & Uncon. \\
\midrule
0.10 & &
  $11.9{\scriptstyle\pm3.7}$ & $13.3{\scriptstyle\pm2.6}$ &
  $14.2{\scriptstyle\pm3.0}$ & $11.9{\scriptstyle\pm0.9}$ &
  $\approx$100 & $\approx$100 & $\approx$100 & $\approx$100 &
  $13.5{\scriptstyle\pm5.7}$ & $10.8{\scriptstyle\pm0.5}$ &
  $11.8{\scriptstyle\pm1.3}$ & $11.2{\scriptstyle\pm1.8}$ \\
0.20 & &
  $22.5{\scriptstyle\pm2.6}$ & $49.5{\scriptstyle\pm23.6}$ &
  $13.1{\scriptstyle\pm1.3}$ & $25.1{\scriptstyle\pm11.3}$ &
  $\approx$100 & $\approx$100 & $\approx$100 & $\approx$100 &
  $14.8{\scriptstyle\pm2.2}$ & $17.8{\scriptstyle\pm2.8}$ &
  $14.8{\scriptstyle\pm6.2}$ & $17.8{\scriptstyle\pm3.0}$ \\
0.40 & &
  $70.0{\scriptstyle\pm28.2}$ & $66.0{\scriptstyle\pm20.4}$ &
  $29.6{\scriptstyle\pm10.5}$ & $15.7{\scriptstyle\pm2.0}$ &
  $\approx$100 & $\approx$100 & $\approx$100 & $\approx$100 &
  $36.2{\scriptstyle\pm35.8}$ & $36.6{\scriptstyle\pm21.3}$ &
  $42.9{\scriptstyle\pm30.5}$ & $45.6{\scriptstyle\pm16.1}$ \\
\bottomrule
\end{tabular}
\end{table*}

\subsection{Epsilon Calibration}

\begin{table}[t]
\centering
\caption{Perturbation magnitude in original feature units.
         A perturbation of $\varepsilon$ in normalized space corresponds
         to $\varepsilon \times \sigma_i$ in original units.
         Top: UNSW-NB15 representative features.
         Bottom: CICIDS-2017 representative features (duration in
         microseconds per CICFlowMeter convention).}
\label{tab:calibration}
\small
\begin{tabular}{l r r r}
\toprule
\textbf{Feature} & \textbf{$\varepsilon=0.10$} & \textbf{$\varepsilon=0.20$} &
  \textbf{$\varepsilon=0.40$} \\
\midrule
\multicolumn{4}{l}{\emph{UNSW-NB15}} \\
\texttt{dur} (s)    & 0.35 & 0.70 & 1.39 \\
\texttt{sbytes} (B) & 5,245 & 10,489 & 20,978 \\
\texttt{spkts}      & 7.46 & 14.91 & 29.82 \\
\texttt{dbytes} (B) & 16,187 & 32,373 & 64,746 \\
\texttt{dpkts}      & 12.20 & 24.41 & 48.81 \\
\texttt{sttl}       & 7.46 & 14.92 & 29.83 \\
\texttt{dttl}       & 4.28 & 8.57 & 17.13 \\
\midrule
\multicolumn{4}{l}{\emph{CICIDS-2017}} \\
\texttt{Flow Duration} ($\mu$s) & 3,368,619 & 6,737,239 & 13,474,477 \\
\texttt{Total Fwd Packets}      & 53.3 & 106.6 & 213.3 \\
\texttt{Total Bwd Packets}      & 71.2 & 142.4 & 284.9 \\
\texttt{Fwd Packet Length Max} (B) & 72.4 & 144.8 & 289.5 \\
\texttt{Flow IAT Mean} ($\mu$s) & 454,481 & 908,962 & 1,817,924 \\
\bottomrule
\end{tabular}
\end{table}

\subsection{Discussion}
\label{sec:discussion}

\textbf{Constraint tightness governs gap magnitude.}
The key structural finding is that the constraint inflation gap scales
with how many features have tight documented bounds.
UNSW-NB15 and NSL-KDD have 15--20 features constrained to narrow
ranges (TTL, boolean flags, rate features strictly in $[0,1]$, count
ceilings).  Unconstrained PGD exploits all of these simultaneously,
producing evasion rates of 79--84\% at the same $\varepsilon$ that
constrained PGD achieves 7--17\%.
CICIDS-2017 has 78 features but most carry only a non-negativity
lower bound; the three features with no lower bound (the IAT min
fields) absorb further gradient mass in unconstrained attacks.
The result is a smaller but consistently positive gap throughout the
$\varepsilon$ sweep.
This relationship---tighter constraints, larger inflation---has a direct
interpretation: every percentage point of gap represents evasion
\emph{purchased by violating domain physics}.  Prior evaluations that
omit projection are not measuring adversarial robustness; they are
measuring how easily an optimizer can find physically impossible inputs.

\textbf{Transfer is architecture- and dataset-dependent.}
On CICIDS-2017, gradient-based attacks from the MLP surrogate
transfer at $\approx\!100\%$ to both RF and XGBoost regardless of
whether domain constraints are applied.
This indicates that all three classifiers have nearly identical decision
boundaries on this feature space, likely because the 78 all-numeric
CICFlowMeter features provide a smooth, aligned representation across
architectures.
On UNSW-NB15 and NSL-KDD, transfer rates are moderate (10--70\%) and
highly variable (standard deviations up to 35 pp at $\varepsilon=0.40$),
consistent with the mixed-type feature space (numeric + one-hot
categorical) creating architecture-specific decision boundaries that
gradient surrogates do not generalize across.
Crucially, \emph{the constraint inflation gap in white-box evasion is
not corrected by this transfer variability}: even on datasets where
constrained transfer exceeds unconstrained transfer (e.g., UNSW-NB15 at
$\varepsilon=0.40$, RF), the \emph{absolute} constrained transfer rates
remain well below unconstrained white-box numbers.

\textbf{Adaptive adversary upper bound.}
As argued in Section~\ref{sec:adaptive} and proved in
Appendix~\ref{app:adaptive}, constrained PGD is the optimal
adaptive attack within the $L_\infty +$ ConstraintSpec threat model.
Our constrained evasion rates therefore represent the ceiling of what a
domain-aware adversary can achieve against these classifiers at these
budgets---not a lower bound that a cleverer attack could exceed.

\textbf{Error bars confirm directional robustness.}
Across all 3 seeds, the constraint inflation gap (unconstrained PGD
$-$ constrained PGD) is positive in 52 of 54 dataset-$\varepsilon$
configurations ($96.3\%$).  The two negative gaps occur only at
$\varepsilon = 0.50$ on UNSW-NB15, where both constrained and
unconstrained attacks achieve near-ceiling evasion ($>99\%$) and
random variation dominates.  The directional claim is robust to
seed choice.

% =============================================================
\section{The \texttt{netadv} Library (v0.2)}
\label{sec:tool}

The \texttt{netadv} library provides a reusable, dataset-agnostic
implementation of domain-constrained adversarial evaluation for NIDS.
Version 0.2 (this paper) extends the AISec release with the following
modules:

\begin{itemize}
  \item \texttt{classifiers/sklearn\_wrap.py} — \texttt{SklearnClassifier}
        wrapper implementing a common \texttt{predict}/\texttt{predict\_proba}
        interface for RF and XGBoost; used by the transfer evaluation loop.
  \item \texttt{attacks/transfer.py} — \texttt{transfer\_evasion\_rate()}
        and \texttt{transfer\_table()}: evaluate pre-generated adversarial
        examples against a list of black-box target classifiers.
  \item \texttt{eval/calibration.py} — \texttt{calibration\_table()}:
        converts $\varepsilon$ from normalized space to original-domain
        units using fitted scaler statistics.
  \item \texttt{constraints/datasets/nsl\_kdd.py} — ConstraintSpec for
        the 38 NSL-KDD numeric features, with rate, count, and boolean
        bounds from the KDD Cup feature documentation.
  \item \texttt{tests/} — 65 unit tests covering constraint projection
        correctness, $L_\infty$ budget enforcement, evasion rate
        denominator semantics, and transfer evaluation logic.
\end{itemize}

% =============================================================
\section{Conclusion}
\label{sec:conclusion}

We have shown that constraint inflation---the overstatement of
adversarial evasion caused by omitting domain constraint projection---is
a systemic property of unconstrained NIDS evaluation, not an artefact
of any single architecture or dataset.
On UNSW-NB15, unconstrained PGD overstates evasion by \textbf{72.0 pp}
at $\varepsilon=0.20$; on NSL-KDD by \textbf{67.4 pp} at $\varepsilon=0.10$.
The gap is smaller on CICIDS-2017 (up to 12.2 pp), where looser feature
semantics limit the volume of infeasible space available to the optimizer.
This relationship between constraint tightness and gap magnitude is itself
a finding: it provides a mechanism---not just an observation---and a
prediction that future evaluations can test against new datasets.

Transfer results reveal a further dataset-specific structure: on
CICFlowMeter-based features, surrogate gradients transfer near-perfectly
to tree classifiers; on mixed-type features with categorical components,
transfer is moderate and highly variable.
Neither pattern diminishes the core white-box inflation finding, whose
magnitude far exceeds the transfer uncertainty.

Taken together, these results argue for a methodological standard change
in adversarial NIDS evaluation: domain constraint projection should be
part of every attack implementation, and reported evasion rates should
be accompanied by the ConstraintSpec used.  The \texttt{netadv} library
provides the infrastructure to do this in fewer than ten lines of code.

% =============================================================
\begin{acks}
[Restore after blind review.]
\end{acks}

% =============================================================
\bibliographystyle{ACM-Reference-Format}
\bibliography{catt_ccs}

% =============================================================
\appendix

\section{Adaptive Adversary Proof Sketch}
\label{app:adaptive}

\textbf{Claim.} Constrained PGD (Algorithm~1) is the
strongest $L_\infty$-bounded adaptive adversary against the
ConstraintSpec projection layer.

\textbf{Proof sketch.}  Let $\Pi_\mathcal{P}$ be the ConstraintSpec
projection (element-wise clamp to $[\ell_i, u_i]$) and
$\Pi_\varepsilon$ be the $L_\infty$ ball projection (clamp to
$[\mathbf{x}_0 - \varepsilon, \mathbf{x}_0 + \varepsilon]$).
Both projections are non-expansive and applied in closed form.
At each step the adversary solves:
\[
  \mathbf{x}^{(t+1)} = \Pi_\varepsilon\!\left(\Pi_\mathcal{P}\!\left(
    \mathbf{x}^{(t)} + \alpha \cdot \operatorname{sign}\!\left(
      \nabla_{\mathbf{x}} \mathcal{L}(\theta, \mathbf{x}^{(t)}, y)
    \right)
  \right)\right).
\]
An adversary who knows $\Pi_\mathcal{P}$ and wishes to maximise loss
subject to $\Pi_\varepsilon \circ \Pi_\mathcal{P}$ cannot do better
than taking the gradient step and then projecting, because the gradient
of the \emph{projected} loss
$\mathcal{L}(\theta, \Pi_\mathcal{P}(\mathbf{x}), y)$
with respect to the \emph{unprojected} $\mathbf{x}$ is zero for
features that hit their bounds (the clamp's subgradient is zero at the
boundary).  Constrained PGD is therefore a fixed point of the
best-response problem under this threat model. $\square$

\section{Open Science}
\label{app:science}

All code, constraint specifications, and benchmark scripts are available
at [anonymous repository].  The repository contains:
\begin{itemize}
  \item \texttt{netadv/} — core library (constraints, attacks, evaluation).
  \item \texttt{benchmarks/unsw\_nb15/run\_ccs.py} — reproduces
        Tables~\ref{tab:unsw_wb}, \ref{tab:transfer},
        \ref{tab:calibration}.
  \item \texttt{benchmarks/cicids2017/run\_ccs.py} — reproduces
        Tables~\ref{tab:cicids_wb}, \ref{tab:transfer}.
  \item \texttt{benchmarks/nsl\_kdd/run\_ccs.py} — reproduces
        Tables~\ref{tab:nsl_wb}, \ref{tab:transfer}.
\end{itemize}
UNSW-NB15 and CICIDS-2017 are publicly available from their respective
authors.  NSL-KDD is available at the University of New Brunswick
dataset repository.

% =============================================================
% Generative AI Disclosure
% =============================================================
\section*{Generative AI Disclosure}

Claude Code (Anthropic, claude-sonnet-4-6) was used during the conduct
of this research to assist with software implementation, including the
\texttt{netadv} library extensions (transfer attack module, sklearn
classifier wrappers, epsilon calibration utilities, NSL-KDD constraint
specification, Colab benchmark notebooks), unit test suite, and as a
writing assistant during manuscript preparation.  All experimental
design decisions, domain knowledge, constraint specification choices,
result interpretation, and accountability for this paper's claims rest
with the authors.

\end{document}
